\documentclass[a4paper,twoside]{article}

\usepackage{morewrites}

\usepackage[svgnames,dvipsnames,x11names]{xcolor}
\usepackage{graphicx}
\usepackage[english]{babel}
\usepackage[colorlinks=true, linkcolor=NavyBlue, urlcolor=NavyBlue, citecolor=NavyBlue]{hyperref}
\usepackage[a4paper]{geometry}
\usepackage{ts-fonts}
\usepackage{tcolorbox}
\usepackage{fvextra}
\usepackage{amsmath}
\usepackage{float}
\usepackage{seqsplit}
\usepackage{ts-code}

\usepackage{lastpage}
\usepackage{refcount}
\makeatletter
\def\ps@plain{
\let\@oddhead\@empty
\let\@evenhead\@empty
\def\@oddfoot{
\hfil
\ifnum\getpagerefnumber{LastPage}>1\relax
\thepage
\fi
\hfil}
\let\@evenfoot\@oddfoot
}
\makeatother

\title{Typesetting controls}
\author{}\date{}
\begin{document}
\maketitle
TeXSmith bundles a \tscodeinline{ts-\allowbreak{}typesetting} fragment that tweaks basic paragraph layout, line spacing, and optional line numbers. By default it stays silent---nothing is injected unless you set one of its options.
\section{Configuration}
All keys live under \tscodeinline{press.typesetting} (short aliases also work: \tscodeinline{press.paragraph}, \tscodeinline{paragraph}, etc.).
\begin{tscode}[lang=yaml, engine=pygments]
\begin{Verbatim}[commandchars=\\\{\},breaklines, breakanywhere, commandchars=\\\{\}]
\PY{n+nt}{press}\PY{p}{:}
\PY{+w}{  }\PY{n+nt}{typesetting}\PY{p}{:}
\PY{+w}{    }\PY{n+nt}{paragraph}\PY{p}{:}
\PY{+w}{      }\PY{n+nt}{indent}\PY{p}{:}\PY{+w}{ }\PY{l+lScalar+lScalarPlain}{auto}\PY{+w}{    }\PY{c+c1}{\PYZsh{} true | false | auto}
\PY{+w}{      }\PY{n+nt}{spacing}\PY{p}{:}\PY{+w}{ }\PY{l+lScalar+lScalarPlain}{1cm}\PY{+w}{    }\PY{c+c1}{\PYZsh{} any TeX length; omit to keep the template default}
\PY{+w}{    }\PY{n+nt}{leading}\PY{p}{:}\PY{+w}{ }\PY{l+lScalar+lScalarPlain}{onehalf}\PY{+w}{  }\PY{c+c1}{\PYZsh{} single | onehalf | double | \PYZlt{}length\PYZgt{} | \PYZlt{}number factor\PYZgt{}}
\PY{+w}{    }\PY{n+nt}{lineno}\PY{p}{:}\PY{+w}{ }\PY{l+lScalar+lScalarPlain}{true}\PY{+w}{      }\PY{c+c1}{\PYZsh{} turn on margin line numbers}
\end{Verbatim}
\end{tscode}
\begin{itemize}
\item \tscodeinline{paragraph.indent} controls \tscodeinline{\textbackslash{}parindent} plus the memoir/article \tscodeinline{\textbackslash{}@afterindent…} switches. \tscodeinline{auto} leaves the first paragraph flush and indents the following ones; \tscodeinline{true} always indents; \tscodeinline{false} disables indentation.
\item \tscodeinline{paragraph.spacing} sets \tscodeinline{\textbackslash{}parskip} to your length (leave empty for the template's original value).
\item \tscodeinline{leading} sets line spacing. \tscodeinline{single}, \tscodeinline{onehalf}, and \tscodeinline{double} call the usual spacing commands. A numeric value applies a stretch factor (\tscodeinline{1.2} \tsscript{symbols}{\(\rightarrow\) }1.2\(\times\)). A length sets \tscodeinline{\textbackslash{}baselineskip} directly (\tscodeinline{1em}, \tscodeinline{14pt}, etc.).
\item \tscodeinline{lineno: true} loads the \tscodeinline{lineno} package and enables \tscodeinline{\textbackslash{}linenumbers} for the whole document.
\end{itemize}
\subsection{Class-aware spacing}
\begin{itemize}
\item On \tscodeinline{memoir}, the fragment uses the class-provided \tscodeinline{\textbackslash{}SingleSpacing}, \tscodeinline{\textbackslash{}OnehalfSpacing}, and \tscodeinline{\textbackslash{}DoubleSpacing} when available, falling back to \tscodeinline{\textbackslash{}baselinestretch} updates.
\item On other classes (article, report...), it loads \tscodeinline{setspace} and uses \tscodeinline{\textbackslash{}singlespacing}, \tscodeinline{\textbackslash{}onehalfspacing}, or \tscodeinline{\textbackslash{}doublespacing}.
\end{itemize}
\subsection{Templates}
The built-in \tscodeinline{article} and \tscodeinline{book} templates inject \tscodeinline{ts-\allowbreak{}typesetting} directly (no \tscodeinline{\textbackslash{}usepackage} file to manage). If you don't set any of the options above, the fragment emits nothing and the templates' stock spacing stays unchanged.
\end{document}
